% RecSys Odyssey repository-backed lecture note.
% Add figures with: \coursefigure{figures/name.svg}{Caption}
\section{Multi-head attention}

\begin{abstract}
Several attention subspaces can represent different transition patterns at once.
\end{abstract}

\subsection{Concept}
One attention map must compress every relationship into one pattern. Multi-head attention gives each head separate projections, allowing different heads to specialize in recency, repetition, category transitions or long-range dependencies. Their outputs are concatenated and mixed. More heads do not automatically mean more quality: model width, data, regularization and latency must support the added capacity.

\begin{equation*}
MHA(X) = Concat(head₁, …, headₕ)Wᴼ
\end{equation*}

\subsection{Key terms}
\begin{itemize}
\item \textbf{Head.} One independently projected attention computation.
\item \textbf{Projection.} A learned linear transformation into a head-specific subspace.
\item \textbf{Model width.} The dimensional capacity shared across the Transformer representation.
\end{itemize}
